% cf-ch4b.tex — 译自 FinalVersion.tex L11425-11570
\section{逻辑切分与逻辑乘积}
本节展示如何利用复合（composition）算法 \ref{tractablelayoutcompositionalgorithm} 来计算逻辑切分（logical division）与逻辑乘积（logical product）。
\subsection{逻辑切分的例子}
回顾：若 $A$ 与 $B$ 是布局（layout），则逻辑切分 $A \oslash B$ 定义为
\[
A \oslash B = A \circ (B,B^c)
\]
其中
\[
B^c = \comp(B, \size(A)).
\]

\begin{example}
    设我们想要计算逻辑切分 $A \oslash B$，其中 $A = (8,8):(8,1)$，$B = (2,2):(1,4)$。于是可以把它们写作 $A = L_g$、$B = L_h$ 与 $B^c = L_{h^c}$，其中 $f$ 与 $f^c$ 是如下所示的元组态射（tuple morphism）。
    \[ \begin{tikzcd} [row sep = 1, column sep = 8]
     & & 2 & & & & 2 & & & & \\
     & & 2 & & & & 2 & & & & \\
    2 \ar[urr,mapsto] & & 2 & & 2 \ar[uurr,mapsto] & & 2 & & 8 \ar[drr,mapsto] & & 8\\
    2 \ar[rr,mapsto] & & 2 & & 2 \ar[urr,mapsto] & & 2 & & 8 \ar[urr,mapsto] & & 8\\
     & h & & & & h^c & & & & g &   \\
    \end{tikzcd} \]
    由此可知，$(B,B^c)$ 由如下所示的嵌套元组态射（nested tuple morphism）$f = (h,h^c)$ 编码。
    \[ \begin{tikzcd} [row sep = 1, column sep = 8]
      & & 2 \ar[rr,mapsto] & & 2 \\
      & & 2 \ar[drr,mapsto] & & 2\\
    4 \ar[urr,-] \ar[uurr,-] & & 2 \ar[urr,mapsto] & & 2 \\
    4 \ar[rr,-] \ar[urr,-]  & & 2 \ar[rr,mapsto] & & 2\\
    & & & f &
    \end{tikzcd} \]
    接下来，我们像之前那样执行复合算法。我们用算法 \ref{mutualrefinementalgorithm} 求出相互加细（mutual refinement）
    \[ \begin{tikzcd} [row sep = 1, column sep = 8]
       & & 4 & & \\
       2 \ar[rr,-] & & 2 & & \\
       2 \ar[rr,-]& & 2 & & \\
     2 \ar[rr,-] & & 2 & & \ar[uull,-] \ar[uuull,-] 8\\
     2 \ar[rr,-] & & 2 & & \ar[ll,-] \ar[ull,-] \ar[uull,-] 8\\
    \end{tikzcd} \]
    构成图表
    \[ \begin{tikzcd} [row sep = 1, column sep = 8]
      & & & & & & 4 & & & & \\
      & & 2 \ar[rr,mapsto] & & 2 \ar[rr,-] & & 2 & & & & \\
      & & 2 \ar[drr,mapsto] & & 2 \ar[rr,-]& & 2 & & & & \\
    4 \ar[urr,-] \ar[uurr,-] & & 2 \ar[urr,mapsto] & & 2 \ar[rr,-] & & 2 & & \ar[uull,-] \ar[uuull,-] 8 \ar[drr,mapsto] & & 8\\
    4 \ar[rr,-] \ar[urr,-]  & & 2 \ar[rr,mapsto] & & 2 \ar[rr,-] & & 2 & & \ar[ll,-] \ar[ull,-] \ar[uull,-] 8 \ar[urr,mapsto] & & 8\\
    & & & f & & & & & &  g &
    \end{tikzcd} \]
    解出该图表，得到
    \[ \begin{tikzcd} [row sep = 1, column sep = 8]
    & & & & 4 \ar[dddrr,mapsto]& & 2 \\
      & & 2 \ar[rr,mapsto] & & 2  \ar[dddrr,mapsto] & & 2\\
      & & 2 \ar[drr,mapsto] & & 2  \ar[uurr, mapsto] & & 2\\
    4 \ar[urr,-] \ar[uurr,-] & & 2 \ar[urr,mapsto] & & 2 \ar[uurr,mapsto] & & 4\\
    4 \ar[rr,-] \ar[urr,-]  & & 2 \ar[rr,mapsto] & & 2  \ar[uurr,mapsto] & & 2\\
    & & & f & & g' &
    \end{tikzcd} \]
    并把 $f$ 与 $g'$ 复合，得到
    \[
    \begin{tikzcd} [row sep = 1, column sep = 8]
      & &   & & 2\\
      & & 2 \ar[dddrr,mapsto] & & 2\\
      & & 2 \ar[urr,mapsto] & & 2\\
    4 \ar[urr,-] \ar[uurr,-] & & 2 \ar[uuurr,mapsto] & & 4\\
    4 \ar[rr,-] \ar[urr,-]  & & 2 \ar[uurr,mapsto] & & 2\\
     & & & g' \circ f & \\
    \end{tikzcd}
    \]
    该嵌套元组态射编码的布局是
    \[
    L_{g' \circ f} = ((2,2),(2,2)):((8,32),(16,1))
    \]
    它在 $((2,2),(2,2))$ 上已合并（coalesced），因此我们断定
    \[
    A \oslash B = ((2,2),(2,2)):((8,32),(16,1)).
    \]
\end{example}

\subsection{逻辑乘积的例子}

回顾：若 $A$ 与 $B$ 是布局，则逻辑乘积 $A \otimes B$ 定义为
\[
A \otimes B = (A, A^c \circ B)
\]
其中
\[
A^c= \comp(A, \size(A)\cdot \cosize(B)).
\]
特别地，若想手工计算 $A \otimes B$，只需算出 $A^c \circ B$，再把所得结果与 $A$ 拼接（concatenate）即可。

\begin{example}
    设我们想要计算逻辑乘积 $A \otimes B$，其中 $A = (2,2):(1,2)$，$B = (5,5):(5,1)$。于是
    \begin{align*}
    A^c & = \comp(A,\size(A)\cdot \cosize(B))\\
     & = \comp(A, 100)\\
     & = (25):(4).
    \end{align*}
    我们按上一节的方式进行。
    \begin{enumerate}
        \item 取 $B$ 与 $\coalesce(A^c) = A^c$ 的标准表示（standard representation）。
        \[ \begin{tikzcd} [row sep = 1, column sep = 8]
        5 \ar[drr,mapsto] & & 5 & & & & 25\\
        5 \ar[urr,mapsto] & & 5 & & 25 \ar[urr,mapsto] & & 4\\
        & f & & & & g & \\
        \end{tikzcd} \]
        \item 应用算法 \ref{mutualrefinementalgorithm}，得到相互加细
        \[ \begin{tikzcd} [row sep = 1, column sep = 8]
        5 \ar[rr,-] & & 5 & & \\
        5 \ar[rr,-] & & 5 & & \ar[ll,-] \ar[ull,-] 25 \\
        \end{tikzcd} \]
        \item 构成图表
        \[ \begin{tikzcd} [row sep = 1, column sep = 8]
        5 \ar[drr,mapsto] & & 5 \ar[rr,-] & & 5 & & & & 25 \\
        5 \ar[urr,mapsto] & & 5 \ar[rr,-] & & 5 & & \ar[ll,-] \ar[ull,-] 25 \ar[urr,mapsto] & & 4\\
         & f & & & & & & g & \\
        \end{tikzcd} \]
        \item 解出该图表，得到
        \[ \begin{tikzcd}[row sep = 1, column sep = 8]
         & & & & 5\\
        5 \ar[drr,mapsto] & & 5 \ar[urr,mapsto] & & 5 \\
        5 \ar[urr,mapsto] & & 5 \ar[urr,mapsto] & & 4\\
        & f' & & g' & \\
        \end{tikzcd} \]
        \item 把 $f'$ 与 $g'$ 复合，得到
        \[ \begin{tikzcd}[row sep = 1, column sep = 8]
         & & 5\\
        5 \ar[rr,mapsto] & & 5\\
        5 \ar[uurr,mapsto] & & 4\\
        & g' \circ f' & \\
        \end{tikzcd} \]
        \item 计算所编码的布局
        \[
        L_{g' \circ f'} = (5,5):(20,4).
        \]
        \item 布局 $L_{g' \circ f'}$ 在 $(5,5)$ 上已合并，因此
        \[
        A^c \circ B = (5,5):(20,4).
        \]
    \end{enumerate}
我们断定
        \[
        A \otimes B = ((2,2),(5,5)):((1,2),(20,4)).
        \]

\end{example}

\appendix

